\scanpage{22}
$E^{2n}$ \deleted{est engendré par}\inserted{admet la \uline{base}} $L_0^n$, \deleted{$E^{-2n}$ engendré par $\Lambda_0^n$, pour \uline{un} choix de $(L,\Lambda)$. Alors pour tt \unsure{$x$}$\in E^{2n}$, on a \unsure{$L=L'$}} \inserted{(pour un $L_0$ de Lefschetz fixé)} \unsure{Or} --- \unsure{où|on}, si $L'$ est un autre élément Lefschetz(ien)
\[
\begin{aligned}
L_0^n&=L_0^n\pi_0=\underset{\rule{.35pt}{.5ex}\rule{3.1em}{.35pt}\rule{.35pt}{.5ex}}{L_0^n{\Lambda'}^{n}}{L'}^{n}
=L_0^n(\beta\Lambda_0^n)(\deleted{\unsure{\beta}}\alpha L_0^n)\\
&=L_0^{\unsure{+n|n}}\alpha\beta\,\pi_0=\alpha\beta\, L_0^n.
\end{aligned}
\]
ce qui prouve que $\alpha\beta=1$, donc $\alpha$ et $\beta$ inversibles, donc \deleted{\unsure{il est}} $L_0^{\unsure{+n|n|\prime n}}$ et $\Lambda_0^{\unsure{+n|n|\prime n}}$ sont des bases de $E^{(2n)}$ et $E^{(-2n)}$ respectivement.

Donc
\[
L\ \unsure{\text{é}|\text{t}|{}^{\prime}}\ \text{Lefschetzien}\quad\Longrightarrow\quad L^n\ \text{\uline{une base de}}\ E^{2n}
\]
\inserted{(et si $\Lambda$\deleted{\unsure{\xi}} associé, $\bar\Lambda$ une base de $E^{-2n}$, \unsure{pour certains $L$})}

\margin{?} On \unsure{sait|saut} \unsure{penser|prouver|poser} une réciproque, dans \deleted{\unsure{le}} \inserted{certains} cas où \deleted{\unsure{$E$}} $E$ est projectif de type fini sur $k$. Ce \deleted{ne} marche certainement pas tel quel, par ex. quand $E=\mathcal L(M^*)$, $M^*$ gradué à degrés $0\leq i\leq4$ ---.

Mais ça ne marche pas non plus pour des anneaux associés \deleted{aux} éléments de degré \unsure{1|2} d'une alg. de Lefschetz, pour la dim 6

[ engendré par $L,L'$ \inserted{de degré\unsure{s|{}} 2}, avec les relations
\[
L^4=0,\ {L'}^{2}-LL'=0,\quad L^2L'-L^3=0,\ L^3L'=0\ ,
\]
donc base
\[
\begin{array}{c@{\ //\ }c@{\ //\ }c@{\ //\ }c}
1&L,L'&L^2,LL'&L^3\\ \hline
\text{degré }0&2&4&
\end{array}
\qquad\Big/\quad
\begin{aligned}
&L\ \text{est lefschetzien,}\\
&L'\ \text{ne l'est pas, bien que}\\
&{L'}^{3}=L^3\quad !\ .
\end{aligned}
\]
% archive sheet number: 21
